\documentclass{article}
\usepackage[utf8]{inputenc}
\usepackage[T1]{fontenc}
\usepackage{geometry}
\usepackage{graphicx}
\usepackage{float}
\usepackage{hyperref}
\usepackage[english,indonesian]{babel}

\title{Optimasi Pemrosesan Template LaTeX Menggunakan Pendekatan Berbasis AI}
\author{Academic Support Agent}
\date{\today}

\begin{document}
\maketitle

\begin{abstract}

\end{abstract}

\chapter{PENDAHULUAN}


\chapter{TINJAUAN PUSTAKA}


% Original content (commented out):
% \documentclass[a4paper]{article}
\usepackage{macros-ohp}
\usepackage{wrapfig}

\renewcommand{\contentsname}{Índice Temático}

\counterwithin{equation}{section}

\renewcommand{\figurename}{Figura}
%\counterwithin{figure}{section}
\renewcommand{\listfigurename}{Índice de Figuras}

\renewcommand{\tablename}{Tabla}
%\counterwithin{table}{section}
%\renewcommand{\listtablename}{Índice de Tablas}

\definecolor{srscolor}{RGB}{255, 255, 255}
\definecolor{srscolor2}{RGB}{199, 221, 141}
\definecolor{blk}{RGB}{0,0,0}
%Please make sure the tex is compiled twice to have all the background images displayed correctly.

\title{\Huge \textbf{ \textcolor{blk}{Mazinger}}}

\author{\Huge \textcolor{blk}{Koji Kabuto}\\[1cm]
	\Large \textcolor{blk}{Asesor(es):}\\[1cm]
	\Large \textcolor{blue}{ Dr. Uno, Dr. Dos }\\
}

\linespread{1.5}

\begin{document}
\begin{titlingpage}
    \tikz[remember picture,overlay] 
    \node[opacity=0.5,inner sep=0pt] at (current page.center){\includegraphics[width=\paperwidth,height=\paperheight]{imgs/Fondo_Protocolo.png}};
\vspace*{3.5cm}
{\let\newpage\relax\maketitle}
\vspace*{\fill}
Zacatecas, Zac.
\end{titlingpage}

\tableofcontents

\newpage
\listoffigures 

\newpage
\listoftables

\newpage

\section{Introducción}

\subsection{Subtema 1 de la introducción}

\subsubsection{Subtema 1 del Subtema 2 de la introducción}

\subsection{Subtema intermedio}

\subsection{Subtema 2 de la introducción}

Recuerda que es preferible que la numeración llegue hasta un 
tercer dígito. Por ejemplo: $3.5.1$. Así, un índice de 
subtema $1.1.1.1$ es preferible evitarlo. 

Recuerda que cuatro dígitos separados por puntos, es para las 
direcciones IP, no para tu tesis. 
Si aún así, usted necesita un apartado más puede hacerlo 
utilizando una notación como la que sigue:

\vspace{1cm}
%\addcontentsline{toc}{section}{\hspace{1cm} Esto no va}
\textbf{Def. Derivada }

Texto que corresponde a la sección que no debería existir.

\newpage
\section{Objetivos}

\subsection{Objetivo General}

\subsection{Objetivos Especificos}

Son un conjunto de pasos que sirven de metas para lograr el 
objetivo general. Si deseas usar una lista numerada, 
podría verse así

\begin{enumerate}
    \item Primer paso
    \item Segundo paso
    \item Tercer paso
    \begin{enumerate}
        \item Paso del paso
        \item El otro paso del paso
    \end{enumerate}
\end{enumerate}

\noindent Si por algún motivo desea hacer una lista 
no numerada, puede verse como sigue:

\begin{itemize}
    \item Primer paso
    \item Segundo paso 
    \item Tercer paso 
    \begin{itemize}
        \item Un subpaso
        \item Otro subpaso
    \end{itemize}
\end{itemize}

\newpage
\section{Marco Teórico}

\subsection{Figuras}
Aquí tenemos un Megaman, como se aprecia en la 
Figura \ref{F_megaman}. Pero tambien podemos ver un 
protoman en la Figura \ref{F_protoman} y para completar 
la triada tenemos un Zero en la Figura \ref{F_zero}. 

\begin{figure}[ht]
    \flushright
    \includegraphics[width=0.3\textwidth]{imgs/megaman.png}
    \caption{Esta es una figura alineada a la derecha.}
    \label{F_megaman}
\end{figure}

\begin{figure}[ht]
    %\flushleft
    \includegraphics[width=0.3\textwidth]{imgs/protoman.png}
    \caption{Esta es una figura de Protoman 
        alineada a la izquierda.}
    \label{F_protoman}
\end{figure}

\begin{figure}[ht]
    \centering
    \includegraphics[width=0.2\textwidth]{imgs/zero.png}
    \caption{Esta es una figura centrada.}
    \label{F_zero}
\end{figure}

\begin{figure}[ht]
    \centering
    %\includegraphics[width=0.3\textwidth]{imgs/megaman.png}
    %\includegraphics[width=0.3\textwidth]{imgs/protoman.png}
    %\includegraphics[width=0.3\textwidth]{imgs/zero.png}
    \includegraphics[height= 4cm]{imgs/megaman.png}
    \includegraphics[height= 4cm]{imgs/protoman.png}
    \includegraphics[height= 4cm]{imgs/zero.png}
    \caption{La fusión de tres imagenes dentro de un solo
    ambiente de figura.}
    \label{F_triada}
\end{figure}

\begin{wrapfigure}{l}{0.25\textwidth}
    \centering
    \includegraphics[width=0.25\textwidth]{imgs/megamanx.png}
    \caption{Evolución a X}
    \label{F_x}
\end{wrapfigure}

\blindtext[5]

\subsection{Ecuaciones}
Para presentar ecuaciones se utiliza el entorno
\textit{equation}.

Para presentar una ecuación numerada, se usa el entorno
tal cual

\begin{equation}
    \label{E_cuadratica}
    16x^2 +3x -12 =0
\end{equation}

Para referenciar lo que pasa en una ecuación, utilizamos su 
etiqueta (label) y usamos el formato que sea preferido por 
el autor. Por ejemplo, podrías decir que la 
ec. (\ref{E_cuadratica}) se resuelve con la famosa 
chicharronera. 

Si deseamos presentar una ecuación que ignore la numeración 
del documento, se usa el entorno \textbf{equation*}

\begin{equation*}
    \label{E_linalg}
    \hat{A} x = \vec{b}
\end{equation*}

pero en este caso, la etiqueta ya no tiene sentido. 
Por ejemplo, la ecuación anterior es la \ref{E_linalg}. Lo 
cual es claramente falso.

\begin{equation}
    \forall \varepsilon \geq 0 \quad 
    \exists \quad \delta \geq 0 \quad \mid 
    |x - a | \leq \varepsilon \Rightarrow  
    |f(x) -L| \leq \delta 
\end{equation}

\newpage
\subsection{Tablas}

Claro que es posible hacer referencia a una tabla dentro 
del texto. Se usa el mismo proceso que para todo lo demás, 
usando el apuntador \textit{ref} hacia la etiqueta adecuada. 
Por ejemplo, podríamos señalar algún resultado ligado con 
la tabla \ref{T_generica}.

\begin{table}[ht]
\centering
\caption{Esta es una tabla genérica.}
\begin{tabular}{|c|c|c|c|c|}
\hline
 T1 & T2  & T3  & T4  & T5  \\ \hline
 a  & b   & c   &  d  & e   \\ \hline
 1  & 2   & 3   &  4  & 5 \\ \hline 
 $\alpha$ & $\beta$  & $\gamma$  & $\delta$  & $\epsilon$  \\ \hline 
\end{tabular}
\label{T_generica}
\end{table}


\section{Citas citables}

Luego \cite{Feynman64} se hizo famoso con la pelicula 
de Oppie.

y con todo lo anterior conseguir hacer levitar una 
rana como en \cite{Simon2000}. Pero también hay una
definición matemática del amor en \cite{Alstola2022}.


\newpage
\bibliographystyle{ieeetr} %ieeetr %unsrt %siam %apalike
\bibliography{refs}
\cite{*}



\end{document}

\end{document}